% TODO: Highlights structure 

\begin{frontmatter}

\begin{abstract}


Growing interest in decarbonization and Arctic accessibility has renewed attention on Europe--Asia shipping corridors. The Northern Sea Route (NSR) is often portrayed as a 30--40\% shortcut relative to Suez, with savings propagated to time, fuel, and CO$_2$. The effect of enforcing sea-only feasibility on
these baselines, and its downstream impact on time, fuel, and CO2, remains
under-examined.

We compare great-circle baselines with sea-only routes computed via A-star search (A*) on a 0.5$^\circ$ grid between Northern Europe and Northeast Asia across the Suez, Cape of Good Hope, and NSR corridors under three waypoint philosophies. Distances are mapped to voyage time using corridor-typical speeds and to fuel/CO$_2$ using main- and auxiliary-engine accounting.

Sea-only routing preserves the ranking NSR $<$ Suez $<$ Cape but compresses NSR’s advantage once realistic speeds are applied. NSR remains shortest (about 8--10k~nm versus 11--12k~nm for Suez), yet typical durations differ modestly and fuel/CO$_2$ savings over Suez are small and variant-dependent. Equal-speed tests restore geometric ordering, and endpoint sensitivity shows larger NSR gains for more northern East Asian ports.

The framework provides a reproducible, corridor-agnostic benchmark for later integration of sea ice, weather, regulatory overlays, and AIS data in dynamic Arctic voyage planning.

\end{abstract}


    % Research Keywords
    \begin{keyword}
    sea-only routing \sep A* shortest-path search \sep great-circle baseline \sep corridor benchmarking \sep Suez Canal \sep Cape of Good Hope \sep Northern Sea Route (NSR) \sep fuel consumption \sep CO$_2$ emissions \sep Arctic shipping
    \end{keyword}


\end{frontmatter}


